% Sections 8.1-8.6, from lectures/L08/appendix/a_state_machines.md (A.1-A.6).

\section{What a finite state machine is}
\label{c8:sec:what}

A finite state machine (FSM) is a circuit whose behaviour depends not only on its current inputs but
on a finite amount of \emph{history}, summarised in a single value called its \textbf{state}. At
every instant the machine is in exactly one state; at every clock edge it may move on to a new one,
determined by its current state and inputs; and its output is computed from that state.

None of this is new hardware. A state machine combines two things you already know: a
\textbf{register} (\chapref{c3:ch}) holding the current state, one flip-flop per state bit, and
\textbf{combinational logic} computing the next state and the output from the current state and the
inputs.

What is new is the discipline. Instead of ad hoc logic feeding an ad hoc set of flip-flops, you work
in a fixed order: first draw a \textbf{state diagram}, a few named states with labelled transitions
between them, and derive the state table only once that is right, then the next-state logic, then
the gate network or the VHDL.

Every synchronous circuit so far is technically a trivial state machine: a counter's state is its
count, a shift register's is its bit pattern. What distinguishes a state machine is that the state
is neither a count nor a shifted pattern but an arbitrary label (\code{STATE_OFF},
\code{STATE_BLINK}, \dots) whose meaning you choose and whose next value may depend on the input
however you like.

There are two varieties, and the difference decides what an output may look at:
\begin{itemize}
  \item \textbf{Moore machines}: the output depends \emph{only} on the current state. That is the
    kind designed by hand here.
  \item \textbf{Mealy machines}: the output depends on the current state \emph{and} the current
    input. \Secref{c8:sec:mealy} covers them briefly, in VHDL only, so that you recognise one when
    you meet it and know the one-clock-cycle difference.
\end{itemize}

% ------------------------------------------------------------------------------------------
\section{Designing a Moore machine by hand}
\label{c8:sec:byhand}

The same process as for every gate network in \chapref{c1:ch} and \ref{c2:ch}, applied to a circuit
with memory. The worked example is a small, closed Moore machine with four states: a circular
counter through a Gray code sequence, stepped one place per button press.

\subsection*{Specification}
\begin{itemize}
  \item Four states, each represented by two flip-flops $\{Q_1, Q_2\}$:

  \begin{narrowtable}{lc}
    \thead{State} & \thead{$\{Q_1, Q_2\}$} \\
    \midrule
    \code{STATE_0} & \code{00} \\
    \code{STATE_1} & \code{01} \\
    \code{STATE_2} & \code{11} \\
    \code{STATE_3} & \code{10} \\
  \end{narrowtable}

  \item A system clock \code{clock}, and an active-low reset \code{reset_n} returning the machine to
    \code{STATE_0}.
  \item An active-low push button \code{button_n}, from which a synchronised, edge-detected pulse
    \code{X} is derived in exactly the way \chapref{c4:ch}'s two-flip-flop synchroniser does.
    \code{X} is high for exactly one clock cycle at a falling edge on \code{button_n}.
  \item An output \code{Y}, connected to an LED, high only in \code{STATE_3}.
  \item The machine is \textbf{closed}: the state after \code{STATE_3} is \code{STATE_0}.
  \item It is a \textbf{Moore} machine: \code{Y} depends only on $\{Q_1, Q_2\}$, never on \code{X}.
\end{itemize}

The encoding is a Gray code ($00 \to 01 \to 11 \to 10 \to 00$), where consecutive states differ in
exactly one bit. That is a deliberate choice, not an accident:
\begin{itemize}
  \item With an ordinary binary count ($00 \to 01 \to 10 \to 11$) a transition such as
    $01 \to 10$ flips two bits at once, and the two flip-flops will not switch at \emph{exactly} the
    same instant.
  \item That skew is harmless to the machine itself. Both flip-flops update at the same clock edge,
    and nothing samples the state register in between, so the machine never acts on the intermediate
    value. It is the same argument as \secref{c5:sec:fmax}'s critical path: as long as everything
    settles before the next edge, what happens in between does not matter.
  \item What the skew \emph{can} disturb is anything reading the state combinationally. An output
    decoded directly from the state bits, like \code{Y} below, can glitch briefly while they are
    skewed. The same goes for a value sampled by \emph{another} clock domain, which is the case
    \secref{c4:sec:sync} warned about.
  \item A Gray code removes both, since only one bit ever changes per transition. It therefore
    matters most for a counter crossing clock domains, and least for a machine whose output is
    registered.
\end{itemize}

Encoding is a design decision in general, and this is the only place in the book where you make it
by hand. Binary uses the fewest flip-flops; Gray minimises the number of bits changing per
transition; \textbf{one-hot} puts one flip-flop per state so that the next-state logic becomes a
single wide OR gate per state, which is often fastest on an FPGA where flip-flops are plentiful. As
soon as you write the machine in VHDL as an enumerated type (\secref{c8:sec:vhdl}), you leave the
decision to the synthesis tool, which picks an encoding for you and re-encodes freely unless you say
otherwise.

\subsection*{State diagram}

\lecfig[0.66]{lectures/L08/appendix/images/fsm_state_diagram.png}{State diagram for the four-state,
Gray-coded Moore machine.}{c8:fig:statediagram}

Every state bubble shows the current state and the output as \code{Q1Q2/Y}, and every arrow is
labelled with the input \code{X} that triggers it. Every state returns to itself as long as
\code{X=0}, and a synchronised pulse from the button (\code{X=1}) steps on to the next state,
wrapping from \code{STATE_3} back to \code{STATE_0}.

\subsection*{State table}
Every combination of current state and input, with the next state $\{Q_1^{+}, Q_2^{+}\}$ and the
output \code{Y}:

\begin{narrowtable}{cccccc}
  \thead{$Q_1$} & \thead{$Q_2$} & \thead{$X$} & \thead{$Q_1^{+}$} & \thead{$Q_2^{+}$} &
    \thead{$Y$} \\
  \midrule
  0 & 0 & 0 & 0 & 0 & 0 \\
  0 & 0 & 1 & 0 & 1 & 0 \\
  0 & 1 & 0 & 0 & 1 & 0 \\
  0 & 1 & 1 & 1 & 1 & 0 \\
  1 & 0 & 0 & 1 & 0 & 1 \\
  1 & 0 & 1 & 0 & 0 & 1 \\
  1 & 1 & 0 & 1 & 1 & 0 \\
  1 & 1 & 1 & 1 & 0 & 0 \\
\end{narrowtable}

\subsection*{Karnaugh-derived logic}
Treating $Q_1^{+}$ and $Q_2^{+}$ separately as Boolean functions of $\{Q_1, Q_2, X\}$ and running
each through a Karnaugh map, the same technique as in \chapref{c2:ch}:

\[ Q_1^{+} = Q_1 X' + Q_2 X \qquad Q_2^{+} = Q_2 X' + Q_1' X \]

Note the shape of both:
\begin{itemize}
  \item each is ``hold when \code{X} is \code{0}, take something from the other bit when \code{X} is
    \code{1}''.
  \item the holding halves, $Q_1 X'$ and $Q_2 X'$, really are mirror images of each other.
  \item the stepping halves are $Q_2 X$ and $Q_1' X$, and the complement on the second is not a
    typo: it is what makes the sequence a Gray code rather than a simple rotation.
\end{itemize}

The output \code{Y} simplifies further, straight from the state table, with no dependence on
\code{X} at all, which confirms that this is a valid Moore machine:

\[ Y = Q_1 Q_2' \]

\subsection*{Realising it as a gate network}
$\{Q_1, Q_2\}$ are one D flip-flop each, exactly like \chapref{c3:ch}'s register: $Q_1^{+}$ and
$Q_2^{+}$ feed the D inputs and $Q_1$, $Q_2$ are read back from the Q outputs. \code{Y} is a single
two-input AND gate with one input inverted, and it reads the current state. That is the whole
machine: two flip-flops and a handful of gates computing $Q_1^{+}$, $Q_2^{+}$ and \code{Y} from
$\{Q_1, Q_2, X\}$.

% ------------------------------------------------------------------------------------------
\section{From gate network to CircuitVerse}
\label{c8:sec:cv}

Before any of this is written in VHDL, realise it by hand in CircuitVerse, exactly as for every gate
network in \chapref{c1:ch}: two-flip-flop synchronise \code{button_n} and \code{reset_n},
edge-detect the synchronised button to produce \code{X}, wire $Q_1^{+}$ and $Q_2^{+}$ to a pair of D
flip-flops, and wire \code{Y} straight into the LED.

\lecfig[1.0]{lectures/L08/appendix/images/fsm_gate_network.png}{Gate network realising the
four-state Moore machine, including two-flip-flop synchronisation and edge detection for
\code{button_n}.}{c8:fig:gatenet}

Read left to right:
\begin{itemize}
  \item \code{button_n} passes through three synchronisation flip-flops
    (\code{button_s1_n}, \code{button_s2_n}, \code{button_s3_n}).
  \item An AND gate derives \code{button_edge_s2} (this circuit's \code{X}) from the second and
    third stages.
  \item The four vertical wires labelled \code{Q1Q2XX'} carry the state bits and the edge pulse on
    to the gates on the right, which implement \secref{c8:sec:byhand}'s equations for $Q_1^{+}$ and
    $Q_2^{+}$.
  \item The two labelled flip-flops store $Q_1$ and $Q_2$ and feed the final AND and NOT gates
    producing \code{Y}, connected to \code{led}.
  \item \code{reset_n} is synchronised by its own pair of flip-flops at the bottom left, exactly
    like every other synchronised reset in this book.
\end{itemize}

Open the runnable circuit yourself: \file{lectures/L08/appendix/images/fsm_gate_network.cv}. Step it
by hand, one button press at a time, and confirm that it visits
\code{STATE_0} $\to$ \code{STATE_1} $\to$ \code{STATE_2} $\to$ \code{STATE_3} $\to$
\code{STATE_0} in order, with the LED lit only in \code{STATE_3}.

% ------------------------------------------------------------------------------------------
\section{The same machine in VHDL}
\label{c8:sec:vhdl}

That was the last time in this book you design a machine by hand. Everything from here writes it
directly in VHDL, which is worth having earned rather than having started with: the synthesis tool
does exactly what you just did on paper, and when a state machine misbehaves it is the diagram you
debug, not the VHDL.

The translation replaces the explicit state table and the Karnaugh equations with an
\textbf{enumerated type} naming every state, a \textbf{case statement} describing every state's
transitions, and a synthesis tool left to work out the flip-flop encoding and the gate-level logic.

\subsection*{The enumerated state type}
Take a three-state machine controlling an LED: \code{STATE_OFF}, \code{STATE_BLINK} (blinking at a
fixed rate) and \code{STATE_ON}. VHDL declares exactly those three values as a new type:

\begin{vhdlcode}
type state_t is (STATE_OFF, STATE_BLINK, STATE_ON);
signal state: state_t;
...
state <= STATE_ON;
\end{vhdlcode}

\subsection*{The case statement pattern}
The machine is nearly always one synchronous process updating \code{state} at a rising edge via a
\code{case}, one branch per state, with an asynchronous reset returning it to a known starting
state. Assume \code{to_next_state} and \code{to_prev_state} are synchronised, edge-detected pulses
produced in exactly the same way as \code{X} in \secref{c8:sec:byhand} and \ref{c8:sec:cv}:

\begin{vhdlcode}
process(clock, reset_s2_n) is
begin
    if (reset_s2_n = '0') then
        -- On reset, go to STATE_OFF.
        state <= STATE_OFF;
    elsif (rising_edge(clock)) then
        case (state) is
            when STATE_OFF =>
                if (to_next_state = '1') then
                    state <= STATE_BLINK;
                elsif (to_prev_state = '1') then
                    state <= STATE_ON;
                end if;
            -- STATE_BLINK and STATE_ON follow the same shape.
            when others =>
                state <= STATE_OFF;
        end case;
    end if;
end process;
\end{vhdlcode}

\code{case (state) is ... end case;} has to cover every value \code{state_t} can take, and
\code{when others} keeps it exhaustive if you add a state later and forget a branch.

What \code{when others} is \emph{not} is a safety net against a corrupted state register. Every
value \code{state_t} can take is already named above it, so the branch is unreachable in the source
and synthesis normally optimises it away rather than building recovery logic. Once Quartus has
re-encoded the machine, most bit patterns are illegal, and getting out of one takes Quartus's Safe
State Machine setting rather than anything you can write in VHDL.

\subsection*{Worked example: \code{fsm_led}}
\file{fsm_led.vhd} implements exactly that three-state machine, extended to two buttons:
\code{button_n(0)} steps forward, \code{button_n(1)} goes back, the states are closed in a loop, and
\code{STATE_BLINK} toggles the LED every 100 ms.

Three subcomponents are reused, all modules you wrote yourself:
\begin{itemize}
  \item \code{reset_sync}: \chapref{c4:ch}'s reset synchroniser, unchanged
    (Exercise~\ref{c4:ex:resetsync}).
  \item \code{button_sync}: \chapref{c4:ch}'s button synchroniser, also unchanged, but instantiated
    with \code{generic map(2)} so that it handles both buttons at once and produces a two-bit pulse
    vector. That is the payoff from that generic: the file is identical to the one \chapref{c4:ch}
    and \ref{c7:ch} instantiate with \code{1} (Exercise~\ref{c4:ex:buttonsync}).
  \item \code{timer}: \chapref{c7:ch}'s timer, unchanged (Exercise~\ref{c7:ex:timer}).
    \code{fsm_led} declares a generic \code{TIMER_TICK_COUNT} of its own and passes it straight on,
    with the default \code{5_000_000} for 100~ms at \code{50 MHz}, and enables the timer only in
    \code{STATE_BLINK}.
\end{itemize}

None of the three is handed out in \file{fsm_led/}: copy your own in before you build the example.
By this point in the book that is four chapters' own modules composed into a single design, which is
what is worth noticing.

\vhdlfile{lectures/L08/fsm_led/fsm_led.vhd}

Two synchronous processes implement the machine:
\begin{itemize}
  \item \code{STATE_PROCESS} is the \code{case} transition logic above, driven by
    \code{to_next_state} and \code{to_prev_state}, which are derived combinationally from
    \code{button_edge_s2}. Each requires the \emph{other} button's pulse to be \code{'0'} in the
    same cycle, so that pressing both close enough in time for their pulses to land on the same
    clock cycle moves the machine nowhere rather than arbitrarily picking a winner. An ambiguous
    input deserves a defined answer, and ``do nothing'' is the one that cannot surprise you.
  \item \code{LED_PROCESS} drives \code{led_s} from the current state alone: off in
    \code{STATE_OFF}, on in \code{STATE_ON}, toggling at every timeout from the timer in
    \code{STATE_BLINK}. Since it never reads \code{button_n} or \code{button_edge_s2} directly, only
    \code{state}, this is a genuine Moore machine: the LED responds to a press only indirectly,
    through the state the press caused.
\end{itemize}

% ------------------------------------------------------------------------------------------
\section{Mealy machines: the one-clock-cycle difference}
\label{c8:sec:mealy}

Everything so far has been a \textbf{Moore} machine, deliberately: Moore is the safer default and
the form you design by hand. A \textbf{Mealy} machine relaxes the rule defining a Moore machine: the
output may depend on the current state \emph{and} the current input, computed combinationally. It is
not stored in a register, and it changes within the same clock cycle as the input does.

\subsection*{The behaviour: detecting two \code{1} bits in a row}
A serial input \code{din} carries one bit per clock cycle, and \code{y} should indicate that the two
most recent bits were both \code{1}, with overlapping hits counted, so \code{111} reports a hit
twice.

As a Mealy machine this needs only two states, since the output can react to the input \emph{before}
it shows up in the next state: \code{STATE_IDLE} (the previous bit was \code{0}, or this is the
first) and \code{STATE_ONE} (the previous bit was \code{1}).

\begin{narrowtable}{lccc}
  \thead{State} & \thead{\code{din}} & \thead{Next state} & \thead{\code{y}} \\
  \midrule
  \code{STATE_IDLE} & 0 & \code{STATE_IDLE} & 0 \\
  \code{STATE_IDLE} & 1 & \code{STATE_ONE}  & 0 \\
  \code{STATE_ONE}  & 0 & \code{STATE_IDLE} & 0 \\
  \code{STATE_ONE}  & 1 & \code{STATE_ONE}  & \textbf{1} \\
\end{narrowtable}

The state process is \secref{c8:sec:vhdl}'s \code{case} pattern; only the output differs:

\vhdlfile{lectures/L08/seq_detect_mealy/seq_detect_mealy.vhd}

The output assignment is a single concurrent statement, entirely outside any clocked process:

\begin{vhdlcode}
y <= '1' when (state = STATE_ONE and din = '1') else '0';
\end{vhdlcode}

That is a \textbf{conditional signal assignment}, and it is the last bit of VHDL syntax this book
introduces. It is a concurrent statement, so it belongs in the architecture body rather than inside a
process, and it describes a multiplexer: take the first value whose condition holds, otherwise the
value after \code{else}. It is the same hardware you would get from an \code{if}/\code{else} inside a
combinational process, written in one line instead of five, and it is worth having because a
one-line output assignment makes it obvious at a glance what the output depends on.

\begin{vhdlcode}
-- These two describe exactly the same multiplexer.
x <= a when sel = '1' else b;

process(sel, a, b) is
begin
    if (sel = '1') then
        x <= a;
    else
        x <= b;
    end if;
end process;
\end{vhdlcode}

Always give it an \code{else}. Omit one and you have described a signal keeping its old value when
no condition holds, which is a latch, exactly as in \secref{c5:sec:latch}. \code{fsm_led} uses this
form three times, and Exercise~\ref{c8:ex:mealy101} asks you to use it.

\textbf{That single line is the whole difference.} It reads \code{din}, so it is Mealy. Move the same
condition into the clocked process, or take \code{din} out of it, and you have a Moore machine
instead.

\subsection*{What it costs and what it buys}
The Moore version of this detector needs a \textbf{third} state, \code{M_TWO}, existing purely to
remember for one cycle that the hit occurred, so that the output has something to read that does not
depend on the input. Feed both machines \code{1, 1, 0, 1, 1, 1}:

\begin{narrowtable}{cclcll}
  \thead{Cycle} & \thead{\code{din}} & \thead{Mealy state (before)} & \thead{Mealy \code{y}} &
    \thead{Moore state (before)} & \thead{Moore \code{y}} \\
  \midrule
  1 & 1 & \code{STATE_IDLE} & 0 & \code{M_IDLE} & 0 \\
  2 & 1 & \code{STATE_ONE}  & \textbf{1} & \code{M_ONE}  & 0 \\
  3 & 0 & \code{STATE_ONE}  & 0 & \code{M_TWO}  & \textbf{1} \\
  4 & 1 & \code{STATE_IDLE} & 0 & \code{M_IDLE} & 0 \\
  5 & 1 & \code{STATE_ONE}  & \textbf{1} & \code{M_ONE}  & 0 \\
  6 & 1 & \code{STATE_ONE}  & \textbf{1} & \code{M_TWO}  & \textbf{1} \\
\end{narrowtable}

Moore's \code{y} column is exactly Mealy's \code{y} column shifted down one row. That is the whole
trade-off: Mealy reports the hit in the \emph{same} cycle the second \code{1} arrives where Moore
takes one more, and Mealy needs 2 states where Moore needs 3.

\subsection*{Which to choose}
Moore, nearly always, and that is what the rest of the book uses, the group project included:
\begin{itemize}
  \item its output does not depend on any primary input, so it cannot follow a noisy one.
  \item a Mealy output can in principle glitch combinationally if its inputs are not already clean.
  \item note that ``Moore'' does not by itself mean glitch-free. An output decoded combinationally
    from several state bits, like \secref{c8:sec:byhand}'s $Y = Q_1 Q_2'$, can still glitch while
    those bits are skewed at a transition. What makes an output glitch-free is \emph{registering}
    it, as \code{fsm_led}'s \code{LED_PROCESS} does, and that is what you want on anything driving a
    wire out of the circuit. A serial \code{tx} line decoded straight from the state bits will send
    out spurious edges.
  \item a machine whose output depends on fewer things is a machine with fewer ways to be wrong.
\end{itemize}

Reach for Mealy only when you specifically need the lower latency or the smaller state count, and be
aware that you are then trading a synchronous output for a combinational one.

% ------------------------------------------------------------------------------------------
\section{On the board}
\label{c8:sec:board}

\code{fsm_led} is synthesised and run exactly as every VHDL design has been since \chapref{c1:ch}:
create a Quartus Prime Lite project with the DE0-CV's \code{5CEBA4F23C7N} as the target device, add
the \file{.vhd} files for the example and its subcomponents, compile, assign every port a physical
pin in the Pin Planner, then recompile and program the board.

For \code{fsm_led}, assign \code{clock} the 50 MHz oscillator's pin, \code{reset_n} and
\code{button_n(1 downto 0)} three push buttons, and \code{led} an LED on the board. Then confirm that
repeated presses of ``next'' make the LED cycle off $\to$ blinking $\to$ on $\to$ off, that
``previous'' goes the other way, and that reset returns it to off.

\code{seq_detect_mealy} is deliberately \textbf{not} demonstrated on the board. It consumes one bit
per clock cycle, and a switch flipped by hand holds its value for millions of cycles, so all an LED
can show is the degenerate case: \code{din} held high, \code{y} high from the second cycle on. The
one-clock-cycle difference the whole of \secref{c8:sec:mealy} is about happens far too fast to see.
Run its testbench instead, which is the general lesson rather than an exception:

\begin{shell}
cd lectures/L08/seq_detect_mealy
cp ../../L04/exercises/reset_sync/reset_sync.vhd .     # the one you wrote in L04
ghdl -a --std=93 reset_sync.vhd seq_detect_mealy.vhd seq_detect_mealy_tb.vhd
ghdl -e --std=93 seq_detect_mealy_tb
ghdl -r --std=93 seq_detect_mealy_tb --assert-level=error --stop-time=10ms
\end{shell}

Exact pin numbers depend on which buttons and LEDs on the DE0-CV you pick, assigned the same way as
for every design since \chapref{c1:ch}. Look up the DE0-CV's own documentation for its physical pin
configuration if you need to confirm which header pin corresponds to which labelled button, switch
or LED.
